\chapter{Architektura i struktura aplikacji}
\label{ch:architektura}

\section{Og\'{o}lna struktura warstwowa}

Aplikacja podzielona jest na trzy logiczne warstwy zgodnie ze wzorcem
Model--View--ViewModel~(MVVM):

\begin{itemize}
  \item \textbf{Warstwa modelu} -- klasy encji Entity~Framework mapowane na tabele bazy
        danych (\texttt{Models/}): \texttt{Kursant}, \texttt{Instruktor}, \texttt{Pojazd},
        \texttt{Administrator}, \texttt{RezerwacjaJazdy}.
  \item \textbf{Warstwa widoku} -- pliki XAML i~ich~\emph{code-behind} (\texttt{Views/}):
        okna i~kontrolki u\.{z}ytkownika.
  \item \textbf{Warstwa ViewModel} -- klasy dziedzicz\k{a}ce po \texttt{ViewModelBase}
        (\texttt{ViewModels/}): logika biznesowa, walidacja, komunikacja z~baz\k{a}.
\end{itemize}

Komunikacja pomi\k{e}dzy widokiem a~ViewModel odbywa si\k{e} przez mechanizm
\textbf{wi\k{a}za\'{n} danych}~(data binding) WPF oraz wzorzec polece\'{n}
\texttt{ICommand} zrealizowany klas\k{a} \texttt{RelayCommand}.

\section{Diagram pakiet\'{o}w}

\begin{table}[H]
\centering
\caption{Struktura katalog\'{o}w projektu}
\label{tab:struktura}
\begin{tabular}{|l|l|}
\hline
\textbf{Katalog} & \textbf{Zawarto\'{s}\'{c}} \\
\hline
\texttt{Models/}      & Klasy encji (5 plik\'{o}w) \\
\texttt{Data/}        & \texttt{OskDbContext} -- kontekst EF Core \\
\texttt{ViewModels/}  & 5 klas ViewModel + \texttt{ViewModelBase} + \texttt{RelayCommand} \\
\texttt{Views/}       & Okna~XAML i~code-behind (7 plik\'{o}w) \\
\texttt{Migrations/}  & Migracje EF Core (historia schematu bazy) \\
\hline
\end{tabular}
\end{table}

\section{Przep\l{}yw sterowania}

Po uruchomieniu aplikacji wy\'{s}wietlane jest okno logowania.
Po pomy\'{s}lnym uwierzytelnieniu otwierane jest g\l{}\'{o}wne okno aplikacji,
kt\'{o}re w~zale\.{z}no\'{s}ci od roli u\.{z}ytkownika wy\'{s}wietla odpowiednie
opcje menu i~wczytuje odpowiedni widok do centralnego obszaru.

Klikni\k{e}cie przycisku menu podmienia zawarto\'{s}\'{c} obszaru roboczego na jeden
z~widok\'{o}w: kursanci, instruktorzy, pojazdy lub rezerwacje.
Ka\.{z}dy widok \l{}aduje niezb\k{e}dne dane z~bazy przy swojej inicjalizacji.

\section{Wzorzec MVVM w projekcie}

Zastosowanie wzorca MVVM pozwala na wyra\'{z}ne rozdzielenie interfejsu graficznego
od logiki aplikacji. Dzi\k{e}ki temu zmiany w~wygl\k{a}dzie okna nie wp\l{}ywaj\k{a}
na logik\k{e} dzia\l{}ania programu i~odwrotnie. Przyciski w~oknach s\k{a}
bezpo\'{s}rednio powi\k{a}zane z~odpowiednimi akcjami w~kodzie, bez
dodatkowego kodu po\'{s}rednicz\k{a}cego.

\section{Diagram klas}

Poni\.{z}szy diagram przedstawia hierarchi\k{e} dziedziczenia w~projekcie.
Klasa \texttt{ViewModelBase} jest klas\k{a} bazow\k{a}, po kt\'{o}rej
dziedzicz\k{a} wszystkie ViewModele. Klasa \texttt{RelayCommand} realizuje
interfejs \texttt{ICommand}.

\begin{figure}[H]
  \centering
  \includegraphics[width=\textwidth]{img/diagram_klas}
  \caption{Diagram klas}
  \label{fig:diagram_klas}
\end{figure}

\section{Mechanizmy walidacji danych}

Walidacja danych wej\'{s}ciowych odbywa si\k{e} w~warstwie ViewModel,
przed wykonaniem jakiejkolwiek operacji na bazie danych. Sprawdzane s\k{a}:
\begin{itemize}
  \item poprawno\'{s}\'{c} numeru PESEL (d\l{}ugo\'{s}\'{c}, cyfry, suma kontrolna),
  \item unikalno\'{s}\'{c} loginu i~numeru PESEL w~bazie danych,
  \item wype\l{}nienie wszystkich wymaganych p\'{o}l formularza,
  \item warto\'{s}\'{c} daty rezerwacji -- termin nie mo\.{z}e by\'{c} z~przesz\l{}o\'{s}ci.
\end{itemize}

W~przypadku b\l{}\k{e}du u\.{z}ytkownik otrzymuje czytelny komunikat opisuj\k{a}cy
przyczyn\k{e} problemu, a~operacja nie jest wykonywana.

\section{Obs\l{}uga b\l{}\k{e}d\'{o}w i~wyj\k{a}tk\'{o}w}

Operacje bazodanowe s\k{a} opakowane w~bloki \texttt{try-catch}, dzi\k{e}ki czemu
wyj\k{a}tki (np. b\l{}\k{a}d po\l{}\k{a}czenia z~baz\k{a} lub naruszenie
klucza obcego) s\k{a} przechwytywane i~prezentowane u\.{z}ytkownikowi
w~postaci okna dialogowego z~opisem problemu, zamiast powodowa\'{c}
awari\k{e} aplikacji.